\section{An explicit finite-prime form of the bounded-degree gate}
\label{sec:finite-prime-gate}

The prime distribution input controls the search for a prime outside a
\emph{known} exceptional integer. The geometric degree bound and that
integer remain separate inputs. This distinction makes the consequence
of the cited Dirichlet GRH both finite and testable.

\begin{proposition}[R30]
\label{prop:R30}
Assume Dirichlet GRH, as supplied in \cite{Alpoge2026}. For every positive
integer $b$ and either $a=1$ or $a=2$, there is a rational prime $p$ with
\[
 p\equiv a\pmod3,\qquad p\nmid b,\qquad
 p\leq X(b):=\max\{\exp(14),\,4\log b+4\}.
\]
\end{proposition}

\begin{proof}
Write $\vartheta(x;3,a)=\sum_{p\leq x,\ p\equiv a\ (3)}\log p$.
The theta estimate in \cite[Corollary 1.1 and Table 1, row
$\log x_0=10$, p.~2]{Lee2023} gives, for $x\geq\exp(10)$,
\[
 \left|\vartheta(x;3,a)-\frac{x}{2}\right|
 <\left(\frac{\log x}{8\pi}+\frac{\log3}{2\pi}+9.29179\right)
       \sqrt{x}\log x-8311.79.
\]
Discard the negative constant. For $y=\log x\geq14$, the remaining
error divided by $x$ is less than
\[
 F(y)=\left(\frac{y}{24}+9.5\right)y\exp(-y/2),
\]
using $\pi>3$ and $\log3<1.1$. The derivative has the sign of
\[
 \frac y{12}+9.5-\frac{y^2}{48}-4.75y,
\]
which is negative for $y\geq4$. Also
\[
 4\left(\frac{14}{24}+9.5\right)14
 =\frac{1694}{3}<565<1000
 <\left(\frac{65}{24}\right)^7<\exp(7).
\]
The last inequality follows by truncating the positive series for
$\exp(1)$ after its fourth-degree term. Thus $F(y)<1/4$ and
\[
 \vartheta(x;3,a)>x/4\qquad(x\geq\exp(14)).
\]
If every prime in this congruence class up to $X(b)$ divided $b$, their
product would divide $b$, so that
$\vartheta(X(b);3,a)\leq\log b$. The displayed lower bound instead gives
$\vartheta(X(b);3,a)>X(b)/4\geq\log b+1$, a contradiction.
\end{proof}

\begin{theorem}[R31]
\label{thm:R31}
Fix the arithmetic polarized model and flat class of B04, and fix a
numerator degree bound $D$. Let $b$ be a positive integer divisible by
the exceptional integer $b_D$ of B04 and by every excluded rational
residue characteristic. Suppose a finite set $\mathcal P$ of distinct
rational primes satisfies
\[
 \prod_{p\in\mathcal P}p>b.
\]
For each $p\in\mathcal P$, suppose at least one geometric fiber above
$p$ has effective integral cycles $Z_p^+,Z_p^-$ of degrees at most $D$
with
\[
 [Z_p^+]-[Z_p^-]=N_p t_p,\qquad N_p\in\Z_{>0}.
\]
Then the generic class is algebraic. More precisely, some one of these
representatives lifts with its own multiplier $N_p$ unchanged and with
both degrees at most $D$.

Under Dirichlet GRH it is sufficient to provide these local certificates
for every good rational prime $p\equiv a\pmod3$ up to $X(b)$, for either
fixed choice $a\in\{1,2\}$.
\end{theorem}

\begin{proof}
If every $p\in\mathcal P$ divided $b$, their product would divide $b$,
contrary to the strict inequality. Choose $p\nmid b$. B04 applies to
the indicated fiber above $p$ and lifts its particular pair and
multiplier. The finite-product criterion is unconditional.

For the GRH assertion, R30 supplies a prime in the chosen congruence
class, at most $X(b)$, not dividing $b$. It is good by the definition
of $b$, and its assumed certificate again satisfies B04.
\end{proof}

For the Eisenstein elliptic-product isogeny class, the two congruence
classes in R30 are respectively ordinary and supersingular. The
counterexample in the local-prime section shows why this identification
cannot be imposed on an arbitrary Weil fourfold family from the splitting
behavior of $K$ alone. In the general setting R30 is a statement about
rational congruence classes, with no Newton-stratum assertion.

R31 strengthens the infinitude criterion to a finite certificate once
$b_D$ is known. It does not estimate $b_D$ without the arithmetic
Hilbert-scheme equations, or manufacture a common local degree $D$ from
prime counts. For the fixed T5 lattice and Weil vector, R24 independently
supplies a uniform degree and multiplier bound, and B02 then carries
that bound to Route~2. Neither the uncompleted Hecke statement indexed
under Masullo nor a distribution theorem for arbitrary Newton strata is
needed for that application.
